\section{Proposed Methodology}

\subsection{Scope}

This project builds a pre-deployment diagnostic tool for clinical prediction models. The tool estimates how much a model's performance is likely to drop when used on a new patient population, and recommends whether to deploy it, adjust it, or hold it back. It is important to be clear about what the project does not do: it does not try to fix generalisation failure. Fixing it properly would need causal modelling or domain adaptation techniques, which the literature review shows are not yet practical for most teams. Instead, this project focuses on spotting the risk of failure before it does any harm, and giving practitioners clear guidance based on what is found. The framework works with tabular clinical data that has identifiable subgroups  different hospitals, regions, trial sites, or time periods  and is designed so that any clinical ML team with standard Python skills can use it without needing specialist statistical training.
3.2 Methodological Framework and Architecture Overview
The methodology has three stages that run in order. Each stage answers one question, produces one output, and feeds into the next.
Stage 1 (Divergence Scoring) asks: how different is the new population from the training population? Three metrics are computed between the training and target data. MMD, from Gretton et al. (2012), compares two distributions without assuming any particular shape  useful for clinical data where features are rarely normally distributed. KL divergence, cited by Wörheide et al. (2021) via Duchi and Namkoong (2018), measures how much information is lost when one distribution stands in for another, and is especially good at catching cases where the target population has patients that the training data simply does not cover. Wasserstein distance, cited by Wörheide et al. via Rahimian and Mehrotra (2019), measures the minimum effort needed to reshape one distribution into the other, and handles situations where the two distributions barely overlap more reliably than KL does. Three metrics are used instead of one because Wörheide et al. list all three as standard options for robust optimisation but never actually test which one is most useful for predicting real performance drops  this project directly fills that gap.
Stage 2 (Feature Stability Analysis) asks: which features in the model can be trusted across different settings, and which ones are unreliable? The approach is straightforward: train the model on all available subgroups except one, record how important each feature is using permutation importance, then repeat this for every subgroup. Features that stay consistently important regardless of which subgroup is held out are labelled shift-stable. Features whose importance jumps around are labelled shift-sensitive. This does something similar to what Pavlovic et al. (2024) achieve with causal graphs  finding features whose link to the outcome holds across environments  but without needing anyone to draw a causal diagram first. Cross-context consistency is not the same thing as causal stability, but it can be measured by anyone with basic modelling skills, which directly solves the accessibility problem that makes Pavlovic et al.'s method impractical for most teams.
Stage 3 (Decision Framework) asks: what should the practitioner actually do? It takes the divergence score from Stage 1 and the stability profile from Stage 2, and maps them onto a simple flowchart with four possible outcomes. If divergence is low, deploy the model with routine monitoring. If divergence is high, do not deploy. If divergence is in the middle, look at the feature stability profile: if most key features are stable, deploy with closer monitoring; if stability is mixed, retrain the model using only the stable features; if most features are unstable, restrict the model's use or withdraw it. After deployment, the divergence score is recalculated on incoming data at regular intervals, and if it crosses the upper threshold, the model is flagged for human review. The 0.60 balanced accuracy threshold for acceptable performance comes from Chekroud et al.'s (2024) cross-validation baseline  the best-case internal estimate from their schizophrenia trials.
